% Generated draft: requires mathematical and source review.
\documentclass[11pt]{article}
\usepackage[margin=25mm]{geometry}
\usepackage{fontspec}
\setmainfont{Latin Modern Roman}
\usepackage{amsmath,amssymb,mathtools,unicode-math}
\setmathfont{Latin Modern Math}
\usepackage{enumitem,needspace,xurl,hyperref}
\setlength{\parindent}{0pt}
\setlength{\parskip}{5pt}
\setlength{\emergencystretch}{3em}
\newcommand{\N}{\mathbb{N}}
\newcommand{\Z}{\mathbb{Z}}
\newcommand{\Q}{\mathbb{Q}}
\newcommand{\R}{\mathbb{R}}
\newcommand{\C}{\mathbb{C}}
\newcommand{\F}{\mathbb{F}}
\newcommand{\A}{\mathbb{A}}
\newcommand{\PP}{\mathbb{P}}
\newcommand{\E}{\mathbb{E}}
\newcommand{\eps}{\varepsilon}
\newcommand{\dd}{\,\mathrm{d}}
\begin{document}
\section*{UnsolvedMath ID: 10400097}
\textbf{Source problem code:} catalogue entry
\subsection{Definitions and mathematical statement}
A quandle is a set \(Q\) with an operation \(*\) satisfying \(x*x=x\), bijectivity of each \(S_y:x\mapsto x*y\), and \((x*y)*z=(x*z)*(y*z)\). It is connected when the group generated by the \(S_y\) acts transitively on \(Q\). Isomorphisms are bijections preserving \(*\). The problem is to classify these isomorphism classes for every finite order \(n\geq1\), not merely to give examples or counts.

The dual operation is \(x\mathbin{\bar *}y=S_y^{-1}(x)\); self-dual means isomorphic to this dual. Use \(R_n\) for the dihedral quandle on \(\Z/n\Z\), \(x*y=2y-x\). The source defines \(\Lambda_n=\Z[t^{\pm1}]/(n)\); its module quotients have the usual Alexander operation \(x*y=tx+(1-t)y\). The conjugation quandle convention is \(x*y=y^{-1}xy\).

The supplied historical Table 4 lists the following numbers of connected isomorphism classes, with all lower-bound signs retained:
\[
\begin{array}{c|rrrrrrrrrrrrrrr}
 n&1&2&3&4&5&6&7&8&9&10&11&12&13&14&15\\\hline
 \text{number}&1&0&1&1&3&2&5&\geq3&8&\geq1&9&\geq2&11&\geq0&\geq4
\end{array}
\]
The named entries are:
\begin{itemize}
\item \(n=1\): the trivial quandle; \(n=3\): \(R_3\); \(n=4\): \(\Lambda_2/(t^2+t+1)\).
\item \(n=5\): \(R_5\), \(\Lambda_5/(t-2)\), and its dual.
\item \(n=6\): two subquandles of \(\operatorname{Conj}(S_4)\), not identified more precisely in the excerpt.
\item \(n=7\): \(R_7\), \(\Lambda_7/(t-2)\), \(\Lambda_7/(t-3)\), and the latter two duals.
\item \(n=8\): an abelian extension of \(\Lambda_2/(t^2+t+1)\), and \(\Lambda_2/(t^3+t+1)\) with its dual.
\item \(n=9\): the table names \(R_9\), \(\Lambda_3/(t^2-t+1)\), \(\Lambda_9/(t-2)\), \(\Lambda_3/(t^2+t-1)\), \(R_3\times R_3\), \(\Lambda_3/(t^2+1)\), and dual entries. The flattened self-dual/non-self-dual column alignment does not unambiguously attach every ``their duals'' phrase.
\item \(n=10\): a subquandle of \(\operatorname{Conj}(S_5)\).
\item \(n=11\): \(R_{11}\) and \(\Lambda_{11}/(t-a)\), \(a=2,\ldots,9\).
\item \(n=12\): \(R_3\times(\Lambda_2/(t^2+t+1))\) and an icosahedral quandle.
\item \(n=13\): \(R_{13}\) and \(\Lambda_{13}/(t-a)\), \(a=2,\ldots,11\).
\item \(n=15\): \(R_3\times R_5\), \(R_3\times(\Lambda_5/(t-2))\), its dual, and a subquandle of \(\operatorname{Conj}(S_5)\); the source adds an ellipsis.
\end{itemize}
The prime-order row reports \(p-2\) classes, listing \(R_p\) and \(\Lambda_p/(t-a)\), \(a=2,\ldots,p-2\). Its intended odd-prime range needs checking against the separate order-two row. The icosahedral description uses the vertices and a rotation centered at each vertex; the rotation angle/direction and the extension data above are not specified. These historical entries are not updated or independently certified here.
\subsection{Short English statement}
Classify connected finite quandles up to isomorphism in every order. Preserve the source's historical examples, duality distinctions, and lower-bound counts without treating incomplete table data as a current complete classification.
\subsection{Sources}
\begin{itemize}
\item Local UnsolvedMath registry, permanent ID 10400097, source code catalogue entry; no proposer or proposal year supplied; updated 2026-08-06.
\item Ohtsuki, \emph{Problems on invariants of knots and 3-manifolds} (2002), Problem 5.1, PDF page 84, Table 4: \url{https://msp.org/gtm/2002/04/gtm-2002-04-024s.pdf}. Registry access date 2026-07-29; PDF font-structure extraction. Provenance distinguishes 231 numbered bold headings. Present primary-literature status and redistribution terms require review.
\item The supplied summary mentions small-order partial classifications without exact later references. The table is retained as supplied historical information, not a verified current bound.
\end{itemize}
\end{document}
